%!TeX spellcheck=en_US
\documentclass[11pt,
               a4paper,
               bibtotoc,
               idxtotoc,
               headsepline,
               footsepline,
               footexclude,
               BCOR12mm,
               DIV13,
               openany,   % using this removes blank pages around part / chapter starts.
%               oneside    % include this if you have to print only one page per sheet of paper.
               ]
               {scrbook}

%%% SETTINGS

% no word wrapping
%\righthyphenmin=62
%\lefthyphenmin=62
% fewer hyphens
\usepackage{microtype}

% german symbols
\usepackage[utf8]{inputenc}

% strikethrough by \sout
\usepackage[normalem]{ulem}

% insert graphics
\usepackage{graphicx}
% more flexible figures e.g. graphics with captions beside them
\usepackage{floatrow}
% more flexible captions.
% Use \captionsetup{options} to configure,
% use it in an environment for local setup
\usepackage{caption}
% subfigures (see template):
\usepackage{subcaption}

% more control of enumerations and itemizations
\usepackage{enumitem}
% less space between items
\setlist[itemize]{itemsep=0cm}
\setlist[enumerate]{itemsep=0cm}
% more customizeable tables (e.g. multiple lines per cell)
\usepackage{tabularx}
% fix for vertical centering
\usepackage{ragged2e}
\renewcommand\tabularxcolumn[1]{>{\Centering}m{#1}}
% column types with multiple lines and formatting
\usepackage{array}
\newcolumntype{C}{>{\centering\arraybackslash}X}
\newcolumntype{R}{>{\raggedleft\arraybackslash}X}
\newcolumntype{L}{>{\raggedright\arraybackslash}X}
% merge multiple rows \multirow{2}{*}{bla} & \\ &
\usepackage{multirow}
% activate for tables with page breaking
%\usepackage{ltablex}
% fix for table movement and itemizations
%\keepXColumns

% fix for dynamics spaces after custom commands
\usepackage{xspace}

% tabbing: use with \tab
\usepackage{tabto}
\TabPositions{4cm}

%% fancy math
% propper matrices, underbrace text
%\usepackage{amsmath}
\usepackage{mathtools}
% special symbols e.g. squares
\usepackage{amssymb}

%% plotting
\usepackage{pgfplots}
\usepgfplotslibrary{fillbetween}

%%Settings for code
% code placement right there
\usepackage{float}
% code coloring
\usepackage{xcolor}
% code listing
\usepackage{listings}

% flexible multi column style
\usepackage{multicol}

% graphs
\usepackage{tikz}
\usetikzlibrary{shapes.geometric, arrows}
% define some elements
\tikzstyle{startstop} = [rectangle, rounded corners, minimum width=3cm, minimum height=1cm,text centered, draw=black, fill=blue!30]
\tikzstyle{arrow} = [thick,->,>=stealth]

% Some code highlighting styles you can use with lstlistings
% C++ code style similar to default eclipse
\lstdefinestyle{eclipse-cpp} {
    captionpos=bottom,
    language=C++,
    otherkeywords={final},
    basicstyle=\footnotesize,
    numbers=left,
    numberstyle=\small,
    showstringspaces=false,
    tabsize=2,
    frame=single,
    breaklines=true,
    keywordstyle=\bfseries\color[RGB]{127,0,85},
    identifierstyle=\color[RGB]{0,0,192},
    stringstyle=\color[RGB]{42,0,255},
    commentstyle=\color[RGB]{63,127,95},
}

% If no highlighting is intended
\lstdefinestyle{plain}{
}

% fancy algorithms (see template)
\usepackage[ruled, vlined, linesnumbered]{algorithm2e}
\DontPrintSemicolon
\SetKw{KwBy}{by}
\SetKw{KwAnd}{and}

% clickable links and clickable table of content <3
% Options: links with linebreaks
\PassOptionsToPackage{hyphens}{url}\usepackage[bookmarks=false]{hyperref}
\hypersetup{
    colorlinks,
    citecolor=black,
    filecolor=black,
    linkcolor=black,
    urlcolor=black
}
% Alterations to labels used by \autoref{}: Capitalize everyything
\def\chapterautorefname{Chapter}
\def\sectionautorefname{Section}
\def\subsectionautorefname{Subsection}
\def\algorithmautorefname{Algorithm}
\def\subfigureautorefname{Figure}
% for custon stuff like use:
% \hyperref[custom:foo]{Custom~\ref*{custom:foo}}


\usepackage{lipsum} % for filling pages with stuff

% -------------------------------------------------------------------------------
% --------------------------------- Thesis Info ---------------------------------
% -------------------------------------------------------------------------------

% set title, authors and stuff for the cover
% docytype needs xspace because it is used within text.
\def\doctype{Bachelor's Thesis\xspace}
\def\studyProgram{Informatics}
\def\title{Efficient trajectory modelling for space debris evolution}
% don't try translate every technical term if it would sound off
\def\titleGer{Effiziente Modellierung zur Vorausberechnung von Weltraumschrott Flugbahnen}
\def\author{Oliver Bösing}
% Prof
\def\supervisor{Univ.-Prof. Dr. Hans-Joachim Bungartz}
% PhD Candidate
\def\advisor{Fabio Alexander Gratl, M.Sc.}
\def\date{15.09.2021}

\begin{document}
\frontmatter
% -------------------------------------------------------------------------------
% ---------------------------------- COVERPAGE ----------------------------------
% -------------------------------------------------------------------------------

% correct BCOR - undo at the end !!!
\def\bcorcor{0.15cm}
\addtolength{\hoffset}{\bcorcor}
\thispagestyle{empty}
\vspace{4cm}
\begin{center}
    \includegraphics[width=4cm]{templateStuff/tumlogo.pdf}\\[5mm]
    \huge FAKULTÄT FÜR INFORMATIK\\[5mm]
    \large DER TECHNISCHEN UNIVERSITÄT MÜNCHEN\\[24mm]

    {\Large \doctype in \studyProgram}\\[20mm]
    {\huge\bf \title\par}
    \vspace{15mm}
    {\LARGE  \author}
    \vspace{10mm}
    \begin{figure}[h!]
        \centering
        \includegraphics[width=4cm]{templateStuff/informat.pdf}
   \end{figure}
\end{center}

\cleardoubleemptypage

% -------------------------------------------------------------------------------
% ---------------------------------- TITLEPAGE ----------------------------------
% -------------------------------------------------------------------------------

\def\bcorcor{0.15cm}
\addtolength{\hoffset}{\bcorcor}
\thispagestyle{empty}
\vspace{10mm}
\begin{center}
    \includegraphics[width=4cm]{templateStuff/tumlogo.pdf}\\[5mm]
	\huge FAKULTÄT FÜR INFORMATIK\\[5mm]
	\large DER TECHNISCHEN UNIVERSITÄT MÜNCHEN\\[24mm]
	{\Large \doctype in \studyProgram}\\[20mm]
	{\LARGE\bf \title}\\[10mm]
	{\LARGE\bf \titleGer}\\[10mm]
	\begin{tabular}{ll}
		\Large Author:      	& \Large \author \\[2mm]
		\Large Supervisor:  	& \Large \supervisor\\[2mm]
		\Large Advisor:			& \Large \advisor\\[2mm]
		\Large Date:       		& \Large \date
	\end{tabular}
	\vspace{-1mm}
	\begin{figure}[h!]
		\centering
		\includegraphics[width=4cm]{templateStuff/informat.pdf}
	\end{figure}
\end{center}

% undo BCOR correction
\addtolength{\hoffset}{\bcorcor}
\newpage

% -------------------------------------------------------------------------------
% ---------------------------------- DISCLAIMER ---------------------------------
% -------------------------------------------------------------------------------

\cleardoubleemptypage

\thispagestyle{empty}
\vspace*{0.7\textheight}
\noindent
I confirm that this \MakeLowercase{\doctype} is my own work and I have documented all sources and material used.\\

\vspace{15mm}
\noindent
Munich, \date \hspace{5cm} \author
\cleardoubleemptypage

% -------------------------------------------------------------------------------
% ------------------------------- ACKNOWLEDGEMENTS ------------------------------
% -------------------------------------------------------------------------------

\phantomsection
\addcontentsline{toc}{chapter}{Acknowledgements}
\vspace*{2cm}
\begin{center}
    {\Large \bf Acknowledgements}
\end{center}
\vspace{1cm}

\lipsum[1]

\cleardoublepage

% -------------------------------------------------------------------------------
% ---------------------------------- ABSTRACT -----------------------------------
% -------------------------------------------------------------------------------

\phantomsection
\addcontentsline{toc}{chapter}{Abstract}
\vspace*{2cm}
\begin{center}
    {\Large \bf Abstract}
\end{center}
\vspace{1cm}

\lipsum[2]

\cleardoublepage

\phantomsection
\addcontentsline{toc}{chapter}{Zusammenfassung}
\vspace*{2cm}
\begin{center}
    {\Large \bf Zusammenfassung}
\end{center}
\vspace{1cm}

\lipsum[2]

\cleardoublepage

% -------------------------------------------------------------------------------
% ------------------------------ TABLE OF CONTENTS ------------------------------
% -------------------------------------------------------------------------------

\tableofcontents
\thispagestyle{empty}
\cleardoubleemptypage

% -------------------------------------------------------------------------------
% --------------------------------- MAIN MATTER ---------------------------------
% -------------------------------------------------------------------------------

\mainmatter
\part{Introduction and Background}
\chapter{Introduction}
\label{sec:intro}       % labels can be put almost anywhere and can be referencef from anywhere.
Write some useful intro. Here are tips along the way:

\section{Tips}
\subsection{How to Describe}
% optional: set the spacing between columns
\setlength{\columnsep}{30 pt}
When listing several points you have three basic options:
\begin{multicols}{3}
    \begin{itemize}
        \item itemize
        \item enumerate
        \item description
    \end{itemize}

    \vfill\null
    \columnbreak

    \begin{enumerate}
        \item itemize
        \item enumerate
        \item description
    \end{enumerate}

    \vfill\null
    \columnbreak

    \begin{description}
        \item[itemize] short, unordered
        \item[enumerate] short ordered
        \item[description] listing of descriptions. Also nice for longer ones.
    \end{description}

\end{multicols}


\subsection{How to Quote}

\begin{quote}
    "This is a quote!"
\end{quote}

\begin{itemize}
    \item Citations to a source can be made like this \verb|\cite{gratl17task}| =~\cite{gratl17task}
    \subitem Always join text and the citation with a non-breaking space: \verb|text~\cite{foo}|.
    \item Referencing Sections, Figures, Tables, Formulas: \verb|\autoref{sec:intro}| = \autoref{sec:intro}.
    \item Footnotes for url or further notes: \verb|\footnote{\url{https://www.top500.org}}| = \footnote{\url{https://www.top500.org}}
\end{itemize}

\subsection{How to Math}

Use the align environment for equations especially if you want to align them somehow.

\begin{align}
1 + 1 &\ne 3\\
\left(\dfrac{10}{1}\right) - 9 &= 1
\end{align}

% if you need a pagebreak because figure placement is broken:
\clearpage

\section{Environments}

\subsection{How to Figure}

Anything can also be put in multiple columns.

\begin{multicols}{2} % defines an environment with two columns
    \begin{figure}[H] % [H] for HERE
        \centering
        \includegraphics[width=.9\columnwidth]{figures/scenario_clip_rot.png}
        \caption[Example Figure]{Some Caption. Always also include a source if it wasn't created by you!\\
            \tiny{Source: \cite{gratl17task}}}
        \label{fig:exampleLabel1} % labels always have to be placed after the caption
    \end{figure}

    \columnbreak    % start next column

    \begin{figure}[H]
        \centering
        \begin{tikzpicture}
            \node[anchor=south west,inner sep=0] (image) at (0,0) {\includegraphics[width=.9\columnwidth]{figures/scenario_clip_rot.png}};
            \begin{scope}[x={(image.south east)},y={(image.north west)}]
            \draw[red, thin,rounded corners] (.42,.42) rectangle (.58,.6);
            \end{scope}
        \end{tikzpicture}
        \caption[Figure with tikz]{Figures can be drawn on or completely generated with tikz.}
        \label{fig:exampleLabel2}
    \end{figure}
\end{multicols}

\paragraph{Subfigures}
If grouping of several pictures seems reasonable, think about using subfigures. This often comes in handy with plots.

\begin{figure}[H]
    \centering
    \begin{subfigure}[b]{0.33\textwidth}
        \includegraphics[width=\textwidth]{example-image-a}
        \caption{example-image-a}
        \label{fig:example-image-a}
    \end{subfigure}
    \begin{subfigure}[b]{0.33\textwidth}
        \includegraphics[width=\textwidth]{example-image-b}
        \caption{example-image-b}
        \label{fig:example-image-b}
    \end{subfigure}
    \begin{subfigure}[b]{0.33\textwidth}
        \includegraphics[width=\textwidth]{example-image-c}
        \caption{example-image-c}
        \label{fig:example-image-c}
    \end{subfigure}
    \caption{One caption to describe them all.}
\end{figure}

\subsection{How to Algorithm}

\begin{figure}
\begin{algorithm}[H]

% Define custom keywords
\SetKwFunction{KwNot}{not}
% Define custom Functions
\SetKwFunction{Fissorted}{is\_sorted}
\SetKwFunction{Fbogosort}{bogosort}
\SetKwFunction{Fshuffle}{shuffle}
\SetKwProg{Fn}{Function}{:}{}
\KwIn{\tabto{2cm}data array}
\KwOut{\tabto{2cm} data sorted}
\BlankLine

\tcp{Checks if array is sorted}
\Fn{\Fissorted{data}}{
    \For{i $\leftarrow$ 0 \KwTo data.size() - 1}{
    \label{algo:for}            % labels can also be put in the algorithm
        \If{data[i] $>$ data[i+1]}{
            \Return false
        }
    }
    \Return true
}

\tcp{actual algorithm}
\Fn{\Fbogosort{data}}{
    \While{\KwNot \Fissorted{data}}{
        random.\Fshuffle{data}
    }
}

\caption[Bogosort]{Bogosort}
\label{algo:example}
\end{algorithm}
\caption{some description what is happening}
\end{figure}

\clearpage

\subsection{How to Code}
\begin{lstlisting}[style=eclipse-cpp, caption=General form of a typical runner() function., label=code:runner]
void runner(int type, void *data){
    switch(type)
        case taskType1:
            // do stuff using data
        case taskType2:
            // do other stuff using data
}
\end{lstlisting}

\subsection{How to Table}
\begin{table}[H]
    \begin{tabularx}{\columnwidth}{L | C | R}
        \hline
        \hline
        bla left & bla centered\newline over two lines &  bla right\\
        \hline
        bla left & bla centered & \multirow[c]{2}{\hsize}{cell spanning two rows} \\
        \cline{1-2}
        \multicolumn{2}{c|}{cell spanning two columns} & \\
    \end{tabularx}
    \caption[Some Table]{Fancy table that can contain line breaks and extended cells.}
    \label{tab:example}
\end{table}

\appendix
\part{Appendix}
\chapter{Some more stuff}

For everything that does not really belong in the thesis but is good to mention.

\listoffigures

\listoftables

\bibliographystyle{alpha}
\bibliography{literature}

\end{document}
